\documentclass[11pt]{article}
\usepackage[letterpaper,margin=1in]{geometry}
\usepackage{fontspec}
\setmainfont{DejaVu Serif}
\setsansfont{DejaVu Sans}
\setmonofont{DejaVu Sans Mono}[Scale=0.86]
\usepackage{microtype}
\usepackage{amsmath,amsthm,amssymb,mathtools}
\usepackage{booktabs,array,tabularx,longtable}
\usepackage{enumitem}
\usepackage{xcolor}
\usepackage{hyperref}
\usepackage{cleveref}
\usepackage{tcolorbox}

\definecolor{dkblue}{HTML}{173B57}
\definecolor{dkred}{HTML}{8B1E2D}
\hypersetup{colorlinks=true,linkcolor=dkblue,citecolor=dkblue,urlcolor=dkblue,
  pdftitle={Historical search supplement: Davis--Kahan Proposition 4.4}}

\newtcolorbox{statusbox}{colback=dkblue!4,colframe=dkblue!70!black,boxrule=0.7pt,
  arc=1.5mm,left=3mm,right=3mm,top=2.5mm,bottom=2.5mm}

\newcommand{\Hh}{\mathcal H}
\newcommand{\R}{\mathbb R}
\newcommand{\DK}{Davis--Kahan}

\title{\textbf{Historical search supplement}\\[0.4em]
\large Literature, errata, and archival due diligence for the note\\
\emph{A Counterexample to Davis--Kahan Proposition 4.4}}
\author{Jonathan Crall \and Edward Wang\\[0.5em]
\normalsize with GPT-5.5, GPT-5.6, Opus 4.8, and Fable 5}
\date{22 July 2026}

\begin{document}
\maketitle

This supplement records the literature and archival search behind the companion
note's historical claims.  It is separated out because that note is a
mathematical result, and the search methodology --- databases queried, terms
used, errata procedures checked, archival routes still open --- is due-diligence
record rather than mathematics.  The companion note states only the conclusions;
everything supporting them is here.

\section{Related work and historical positioning}
\label{sec:related}

The present result should be distinguished from the previously known fact that
canonical or direct rotations need not solve every global unitary-completion
problem for every unitarily invariant norm.  The new point asserted here is
narrower: the published Davis--Kahan conclusion fails \emph{inside} its stated
real short-angle regime.  Our present historical assessment is:

\begin{statusbox}
General failures of direct-rotation optimality for Schatten $p<2$ norms were
known by 2001, but the published examples lie outside the Davis--Kahan
short-angle regime.  To our knowledge, Proposition~4.4 has not previously been
observed to fail under its stated hypothesis
$\theta_{\max}\leq\pi/3$.
\end{statusbox}

This is deliberately not an unconditional priority claim.  Bibliographic
absence is difficult to prove, particularly for older, non-digitized material,
private correspondence, and subscriber-only citation databases.  The evidence
nevertheless supports the more precise statement that no widely recognized or
indexed correction of the short-angle claim has been found.

\subsection{The original claim and its advertised repair}

The 1970 paper does not merely assert a bound for the operator norm, the square
norm, or the restriction of the displacement to the source subspace.  It claims
that, in a real Hilbert space and under $\Theta\leq\pi/3$, the direct rotation
minimizes $\lVert I-V\rVert$ for \emph{every} unitarily invariant norm
\cite{DavisKahan1970}.  The trace norm is therefore explicitly included.  The
surrounding discussion contrasts Proposition~4.4 with two preceding examples:
one showing failure for a sufficiently large real angle, and another showing
failure in a complex space.  The reality assumption and the $\pi/3$ bound are
presented as repairing those failures.  There is no hypothesis of a single
principal plane, distinct principal angles, or a competitor preserving the
principal-plane decomposition.

\subsection{Historical treatments of the CS decomposition and direct rotation}

Paige and Wei's historical account of the cosine--sine decomposition reviews
the contributions of Davis and Kahan and explains the relation between the CS
decomposition and the direct rotations of Davis and Kato
\cite{PaigeWei1994}.  It is a natural place where a known error in the geometric
extremal theory might have been mentioned, but it does not flag
Proposition~4.4.  This is not proof that the error was unknown in 1994; it is
supporting evidence that no correction had become part of the standard history.

\subsection{The norm-class divide in Bolotnikov--Li--Rodman}

Bolotnikov, Li, and Rodman studied the minimization problem
\[
 \min\{\lVert Z-I\rVert_\Phi:
       Z\text{ unitary and }Z\mathcal M_j=\mathcal N_j\}
\]
for chains of subspaces and unitarily invariant norms
\cite{BolotnikovLiRodman2001}.  Their formulation includes, as the one-subspace
case, the ambient completion problem considered in Davis--Kahan Section~4.
They introduce the class of \emph{$Q$-norms} and prove canonical minimization
for that class.  For Schatten norms, the $Q$-norm condition holds exactly for
\[
  \lVert\cdot\rVert_p,\qquad 2\leq p\leq\infty,
\]
because $\lVert A\rVert_p^2=\lVert A^*A\rVert_{p/2}$.  This encompasses the
Frobenius norm and the operator norm and parallels the passage from the valid
squared-displacement theorem of Davis--Kahan Proposition~4.3.

The same paper explicitly shows that the canonical result is generally false
for Schatten norms $1\leq p<2$; see its Lemma~3.3.  Thus failure for the trace
norm was known in substance by 2001.  However, the explicit negative example is
obtained by setting the cosine parameter in its CS block to $a=0$.
Geometrically,
\[
 a=\cos\theta=0,
 \qquad \theta=\frac\pi2,
\]
so the example is the orthogonal-subspace endpoint and lies completely outside
the Davis--Kahan range $\theta\leq\pi/3$.  The authors conclude that the
canonical theorem is generally false for non-$Q$-norms, but do not investigate
whether a short-angle hypothesis restores it.

This distinction is essential.  The present note does not claim to discover
that trace-like norms can favor a noncanonical completion.  It supplies an
explicit example showing that the particular $\pi/3$ repair proposed by Davis
and Kahan also fails.

\subsection{The 2006 near miss}

Qiu, Zhang, and Li revisited unitarily invariant geometry on Grassmann spaces
\cite{QiuZhangLi2006}.  In their discussion of extremal properties of the
direct rotation, they repeat three conclusions attributed to Davis--Kahan:
minimal restricted displacement, minimal squared displacement, and minimal
full displacement when the largest canonical angle is at most $\pi/3$.  The
third is precisely the claim refuted here.  They describe the $\pi/3$
restriction as undesirable, but do not question its correctness, and instead
prove the different, globally valid statement that a direct rotation minimizes
a logarithmic or geodesic displacement $\lVert\log Z\rVert_\Phi$.

The paper contains an especially close near miss.  In a majorization example it
compares the angle vectors
\[
 \left(\frac\pi4,\frac\pi4\right)
 \quad\text{and}\quad
 \left(\frac\pi2,0\right)
\]
and observes the reversal
\[
 \sin\frac\pi4+\sin\frac\pi4
   =\sqrt2
   >1
   =\sin\frac\pi2+\sin0.
\]
This is essentially the same ``spread versus concentrate'' mechanism used in
our counterexample: two equal moderate rotations can cost more in an additive
trace-like norm than one larger rotation together with one fixed plane.  The
2006 authors apply the observation to a Grassmann metric and do not connect it
to the ambient completion problem $\min\lVert I-Z\rVert_1$.

The historical evidence is stronger because Chi-Kwong Li coauthored both the
2001 paper establishing general failures for non-$Q$-norms and the 2006 paper
repeating the $\pi/3$ assertion.  The 2006 paper also acknowledges helpful
discussions with Chandler Davis.  The natural interpretation is that the
general norm-class obstruction was understood, while the short-angle claim was
still believed to be correct.  This is evidence rather than proof of what the
authors knew.

\subsection{Later structural and norm-dependent literature}

Dou, Shi, Cui, and Du give general explicit descriptions of intertwining
operators and direct rotations of two orthogonal projections
\cite{DouShiCuiDu2017}.  Their work improves structural existence and
representation results associated with Kato, Avron--Seiler--Simon, and
Davis--Kahan, but it does not report a correction to the Section~4 extremal
statement.

Norm-dependent optimality also appears in orthogonal Procrustes problems.  Cape
studies how an optimizer depends on the chosen norm
\cite{Cape2020}.  Such problems optimize frame alignments of the form
$\inf_Z\lVert X-YZ\rVert$ and reinforce the general lesson that a canonical
optimizer may depend on the norm.  They are nevertheless different from the
ambient orthogonal-extension problem
\[
 \inf\{\lVert I-Z\rVert:Z\mathcal U=\mathcal V\}
\]
and do not identify a Davis--Kahan error.

Books, surveys, and later applications found in this review tend to rely on the
safer neighboring statements: minimization of $(I-Z)P$, minimization of the
squared displacement $(I-Z)^*(I-Z)$ in arbitrary unitarily invariant norms, or
minimality for the operator and Frobenius norms.  These claims remain compatible
with both the $Q$-norm theory and the present counterexample.

\subsection{Errata search and the official publication record}

SIAM's current author guidance states that an accepted author erratum is
published as a separate item and attached to the original article's PDF; a
production correction is likewise appended to the original PDF
\cite{SIAMErrataPolicy}.  We found no erratum, correction, expression of
concern, or attached notice on the official article record for DOI
\texttt{10.1137/0707001}, and targeted searches did not locate a SIAM-recorded
correction of Proposition~4.4.  This supports the narrow statement:

\begin{quote}
There is no SIAM-recorded correction of Proposition~4.4 that we could locate.
\end{quote}

It does not exclude an unpublished observation, correspondence, or a correction
buried in unrelated literature.

\subsection{Search limitations and archival route}

The literature search included the exact $\pi/3$ assertion, ``Proposition 4.4''
with Davis--Kahan, direct rotation combined with trace, nuclear, and Schatten
norms, the angle pattern $(\pi/4,\pi/4)$ versus $(\pi/2,0)$, the exact algebraic
costs in the example, and terms such as \emph{erratum}, \emph{correction},
\emph{false}, and \emph{counterexample}.  The only directly relevant hits were
the general $p<2$ failure in 2001 and the 2006 repetition of the short-angle
claim.

The main unresolved bibliographic limitations are subscriber-only forward
citation lists, MathSciNet and zbMATH reviews that were not exhaustively
exported, dissertations and non-digitized notes, and private correspondence.
The University of Toronto holds the Chandler Davis fonds, comprising 59 boxes
of research material and extensive correspondence, with most files open to
researchers \cite{ChandlerDavisFonds}.  The most promising archival targets are
correspondence and drafts from approximately 1967--1971, exchanges with Kahan,
notes for \emph{The Rotation of Eigenvectors by a Perturbation. III}, and later
correspondence concerning unitary-invariant norms and direct rotations.  An
archival request is a genuinely independent historical route that could alter
the priority assessment.

\subsection{Chronology and confidence assessment}

\begin{center}
\begin{tabularx}{0.97\textwidth}{>{\raggedright\arraybackslash}p{0.17\textwidth}
                                  >{\raggedright\arraybackslash}p{0.29\textwidth}
                                  >{\raggedright\arraybackslash}X}
\toprule
\textbf{Source} & \textbf{What it establishes} & \textbf{Relation to this counterexample}\\
\midrule
Davis--Kahan (1970) & All-UI-norm minimization in real spaces for $\Theta\leq\pi/3$ & The false statement itself.\\
Paige--Wei (1994) & History of the CS decomposition and direct rotation & Does not flag Proposition~4.4.\\
Bolotnikov--Li--Rodman (2001) & Positive theory for $Q$-norms; failure for Schatten $p<2$ & Explicit negative example is at $\theta=\pi/2$, not in the short-angle regime.\\
Qiu--Zhang--Li (2006) & Repeats the $\pi/3$ assertion; proves logarithmic metric optimality & Strong evidence that the precise error was still unnoticed; contains a closely related majorization example.\\
Dou--Shi--Cui--Du (2017) & General formulas for direct rotations and intertwining operators & Structural theory, with no correction of the extremal statement.\\
Cape (2020) & Norm-dependent behavior in a Procrustes problem & Different feasible-set geometry; no Davis--Kahan correction.\\
Present note & Real $\R^4$ counterexample with angles $(\pi/4,\pi/4)$ & Refutes Proposition~4.4 inside its advertised regime.\\
\bottomrule
\end{tabularx}
\end{center}

Our confidence levels are as follows:
\begin{itemize}
\item \textbf{High confidence:} there is no widely recognized or mainstream
  published correction of Proposition~4.4.
\item \textbf{Moderately high confidence:} no indexed paper gives an explicit
  counterexample under $\theta_{\max}\leq\pi/3$.
\item \textbf{Moderate confidence:} the real $\R^4$, $(\pi/4,\pi/4)$,
  trace-norm counterexample is mathematically new.
\item \textbf{Low confidence:} no mathematician ever noticed the error
  privately, in obscure literature, or in correspondence.
\end{itemize}

The strongest responsible novelty language is therefore:

\begin{statusbox}
General failures of direct-rotation optimality for Schatten $p<2$ norms were
previously known, but the published examples lie outside the Davis--Kahan
short-angle regime.  To our knowledge, Proposition~4.4 has not previously been
observed to fail under its stated hypothesis
$\theta_{\max}\leq\pi/3$.
\end{statusbox}

Before submission, the highest-value remaining due diligence is an
institutional MathSciNet/zbMATH/Scopus citation export, an archival request to
the Chandler Davis fonds, and direct outreach to experts or surviving authors.
Further ordinary web searching is unlikely to change the assessment as much as
those three routes.

\begin{thebibliography}{99}

\bibitem{BolotnikovLiRodman2001}
V. Bolotnikov, C.-K. Li, and L. Rodman,
\newblock Minimal distortion problems for classes of unitary matrices,
\newblock \emph{Electronic Journal of Linear Algebra}, 8:26--46, 2001.
\newblock \href{https://doi.org/10.13001/1081-3810.1058}{doi:10.13001/1081-3810.1058}.

\bibitem{Cape2020}
J. Cape,
\newblock Orthogonal Procrustes and norm-dependent optimality,
\newblock \emph{Electronic Journal of Linear Algebra}, 36:158--168, 2020.
\newblock \href{https://doi.org/10.13001/ela.2020.5009}{doi:10.13001/ela.2020.5009}.

\bibitem{DavisKahan1970}
C. Davis and W. M. Kahan,
\newblock The rotation of eigenvectors by a perturbation. III,
\newblock \emph{SIAM Journal on Numerical Analysis}, 7(1):1--46, 1970.
\newblock \href{https://doi.org/10.1137/0707001}{doi:10.1137/0707001}.


\bibitem{PaigeWei1994}
C. C. Paige and M. Wei,
\newblock History and generality of the CS decomposition,
\newblock \emph{Linear Algebra and its Applications}, 208--209:303--326, 1994.
\newblock \href{https://doi.org/10.1016/0024-3795(94)90446-4}{doi:10.1016/0024-3795(94)90446-4}.

\bibitem{DouShiCuiDu2017}
Y.-N. Dou, W.-J. Shi, M.-M. Cui, and H.-K. Du,
\newblock General explicit descriptions for intertwining operators and direct rotations of two orthogonal projections,
\newblock \emph{Linear Algebra and its Applications}, 531:575--591, 2017.
\newblock \href{https://doi.org/10.1016/j.laa.2017.06.036}{doi:10.1016/j.laa.2017.06.036}.

\bibitem{SIAMErrataPolicy}
Society for Industrial and Applied Mathematics,
\newblock Information for journal authors: Errors in published articles and publication of errata,
\newblock accessed 22 July 2026.
\newblock \url{https://epubs.siam.org/journal-authors}.

\bibitem{ChandlerDavisFonds}
University of Toronto Archives and Records Management Services,
\newblock Chandler Davis fonds, UTA 1218, 1941--2022, 59 boxes,
\newblock finding aid accessed 22 July 2026.
\newblock \url{https://discoverarchives.library.utoronto.ca/index.php/chandler-davis-fonds}.

\bibitem{QiuZhangLi2006}
L. Qiu, Y. Zhang, and C.-K. Li,
\newblock Unitarily invariant metrics on the Grassmann space,
\newblock \emph{SIAM Journal on Matrix Analysis and Applications},
27(2):507--531, 2005; published online 2006.
\newblock \href{https://doi.org/10.1137/040607605}{doi:10.1137/040607605}.

\end{thebibliography}

\end{document}
